% Standalone problem document, generated from the adjacent README.md.
% Compile with XeLaTeX. Mathematical content is maintained in Markdown.
\documentclass[11pt,a4paper]{article}
\usepackage[fontsize=10.5pt]{scrextend}
\usepackage[margin=22mm,headheight=15pt,headsep=9mm,footskip=11mm]{geometry}
\usepackage{fontspec}
\setmainfont{texgyrepagella}[Extension=.otf,UprightFont=*-regular,BoldFont=*-bold,ItalicFont=*-italic,BoldItalicFont=*-bolditalic]
\setsansfont{texgyreheros}[Extension=.otf,UprightFont=*-regular,BoldFont=*-bold,ItalicFont=*-italic,BoldItalicFont=*-bolditalic]
\setmonofont{texgyrecursor}[Extension=.otf,UprightFont=*-regular,BoldFont=*-bold,ItalicFont=*-italic,BoldItalicFont=*-bolditalic,Scale=MatchLowercase]
\usepackage{amsmath,amssymb}
\usepackage{unicode-math}
\setmathfont{texgyrepagella-math.otf}
\usepackage{xcolor,microtype,enumitem,needspace,fancyhdr}
\usepackage{bookmark}
\definecolor{ink}{HTML}{17344B}
\definecolor{accent}{HTML}{197D80}
\definecolor{muted}{HTML}{596773}
\hypersetup{colorlinks=true,linkcolor=accent,urlcolor=accent,pdftitle={MF-21 - The
uniform expansion threshold for Toeplitz symbols with higher-order
zeros},pdfauthor={Open Problems in Numerical Linear Algebra}}
\urlstyle{same}
\setlength{\parindent}{0pt}
\setlength{\parskip}{5pt plus 1pt minus 1pt}
\linespread{1.04}
\setlength{\emergencystretch}{3em}
\widowpenalty=10000
\clubpenalty=10000
\displaywidowpenalty=10000
\allowdisplaybreaks[1]
\setlist{nosep,leftmargin=1.5em}
\providecommand{\tightlist}{\setlength{\itemsep}{0pt}\setlength{\parskip}{0pt}}
\providecommand{\pandocbounded}[1]{#1}
\pagestyle{fancy}
\fancyhf{}
\fancyhead[L]{\sffamily\footnotesize\color{muted}OPEN PROBLEMS IN NUMERICAL LINEAR ALGEBRA}
\fancyhead[R]{\sffamily\bfseries\footnotesize\color{accent}MF-21}
\fancyfoot[L]{\parbox[b]{0.9\textwidth}{\sffamily\fontsize{8}{10}\selectfont\color{muted}Editorial ratings; a bounded search cannot certify that a problem remains unsolved.\\Verification check: 2026-09-20}}
\fancyfoot[R]{\sffamily\footnotesize\color{muted}\thepage}
\renewcommand{\headrulewidth}{0.3pt}
\renewcommand{\headrule}{\hbox to\headwidth{\color{accent}\leaders\hrule height \headrulewidth\hfill}}
\makeatletter
\renewcommand{\section}{\Needspace{5\baselineskip}\@startsection{section}{1}{0pt}{12pt plus 2pt minus 2pt}{4pt}{\sffamily\bfseries\large\color{ink}}}
\renewcommand{\subsection}{\Needspace{4\baselineskip}\@startsection{subsection}{2}{0pt}{9pt plus 1pt minus 1pt}{3pt}{\sffamily\bfseries\normalsize\color{ink}}}
\makeatother
\setcounter{secnumdepth}{0}
\begin{document}
{\sffamily\small\color{muted}Matrix functions and stability\par}
\vspace{3pt}
{\raggedright\hyphenpenalty=10000\exhyphenpenalty=10000\sffamily\bfseries\fontsize{21}{24}\selectfont\color{ink}The
uniform expansion threshold for Toeplitz symbols with higher-order
zeros\par}
\vspace{7pt}
{\sffamily\small\textbf{Difficulty:} challenging\quad\textbf{Importance:} interesting
to specialist\par}
{\sffamily\small\bfseries\color{accent}Status: Lean verified\par}
\vspace{4pt}
{\color{accent}\hrule height 0.8pt}
\vspace{6pt}
\textbf{Rating rationale:} Matching bulk asymptotics with extreme
eigenvalues at the exact breakdown order is challenging; the detailed
expansion threshold is mainly important to structured spectral
specialists.

\hypertarget{lean-proof-and-verification-evidence-20-september-2026}{%
\subsection{Lean proof and verification evidence --- 20 September
2026}\label{lean-proof-and-verification-evidence-20-september-2026}}

The
\href{https://github.com/ajt60gaibb/OpenProblemsInNLA/blob/main/matrix-functions-and-stability/MF-21/lean/README.md}{complete
formalization} proves all three original assertions for every \(m\ge3\),
the obstruction for every continuous coefficient family, and the revised
proof's smooth common family with matching upper and finite-head lower
error bounds. The exact Fourier-integral matrix correspondence and
coefficient uniqueness are separately exported. No analytic or
transcendence hypothesis remains unproved in these results.

The immutable local proof commit
\textit{e378ec4679f9119aacf79a0fe30b80493a61ec5b} passed the actual
non-root Linux Comparator, Lean default-kernel replay, standard-axiom
checks, and real sandbox/rejection controls.
\href{https://github.com/ajt60gaibb/OpenProblemsInNLA/blob/main/matrix-functions-and-stability/MF-21/lean/verification/linux-2026-09-20/README.md}{Checked
input snapshot and raw evidence} ·
\href{https://github.com/ajt60gaibb/OpenProblemsInNLA/blob/main/matrix-functions-and-stability/MF-21/lean/Solution.lean}{Five
proved exports} ·
\href{https://github.com/ajt60gaibb/OpenProblemsInNLA/blob/main/matrix-functions-and-stability/MF-21/lean/NUMERICAL_TARGETS.md}{Exact
statement correspondence}.

Separate nonimplementing AI agents approved the
\href{https://github.com/ajt60gaibb/OpenProblemsInNLA/blob/main/matrix-functions-and-stability/MF-21/lean/reviews/independent-fidelity.md}{formal
statement and source correspondence} and
\href{https://github.com/ajt60gaibb/OpenProblemsInNLA/blob/main/matrix-functions-and-stability/MF-21/lean/reviews/independent-proof.md}{proof
and runtime evidence}; an additional
\href{https://github.com/ajt60gaibb/OpenProblemsInNLA/blob/main/matrix-functions-and-stability/MF-21/lean/reviews/independent-cross-review.md}{scoped
cross-review} records authorship exclusions. The formalization and
September proof revisions were produced by OpenAI Codex AI agents.
George Stepaniants retains original manuscript authorship, and the
conjecture and simple-loop method retain the cited authors' attribution.
No external human peer review or source-author endorsement is claimed.
Verification was completed locally before the authorized publication in
\href{https://github.com/ajt60gaibb/OpenProblemsInNLA/pull/309}{PR
\#309}.

\hypertarget{independent-formalization-of-the-frozen-manuscript}{%
\subsection{Independent formalization of the frozen
manuscript}\label{independent-formalization-of-the-frozen-manuscript}}

An
\href{https://github.com/sgstepaniants/OpenProblemsInNLA/tree/fe2183e9b570322d7c1f64bba3480460082af8d3/matrix-functions-and-stability/MF-21/lean-frozen-2026-09-12}{additional
Lean formalization} follows the unchanged manuscript from commit
\textit{eb37bc17}. Its \textit{MF21Restart.Contracts.target\_proved} and
\textit{MF21Restart.Contracts.manuscript\_smooth\_target} prove the
original three-part target with one smooth coefficient family. The
development retains the original boundary determinant, phase indexing,
inverse-kernel limit, and trace contradiction. It supplements the
existing formalization and does not add another distinct solved problem.

Local checks cover all 125 source files and 386 axiom reports, using
only the permitted standard foundational axioms. After the full build,
the final phase-source adjustment was checked by recompiling its 10
affected modules and validating 115 unchanged outputs for reuse. Lean
4.33.1, Mathlib and kernel-mode LeanCert are pinned in the
\href{https://github.com/sgstepaniants/OpenProblemsInNLA/tree/fe2183e9b570322d7c1f64bba3480460082af8d3/matrix-functions-and-stability/MF-21/lean-frozen-2026-09-12}{proof
project}; it includes exact statements, component reviews and two
additional nonimplementing final referees, and the one-process
reproduction command \textit{python3 verify\_local.py}. The
\href{https://github.com/sgstepaniants/OpenProblemsInNLA/actions/runs/35552653565}{separate
Linux run} and
\href{https://github.com/ajt60gaibb/OpenProblemsInNLA/blob/main/references/mf21-frozen-lean-2026-09-21/README.md}{retained
publication evidence} identify the exact checked revision and the actual
Comparator, kernel-replay and sandbox results. George Stepaniants is
credited with his Department of Computing and Mathematical Sciences,
California Institute of Technology affiliation, with substantial AI
assistance disclosed.

\hypertarget{proof-clarification---19-september-2026}{%
\subsection{Proof clarification - 19 September
2026}\label{proof-clarification---19-september-2026}}

The
\href{https://github.com/ajt60gaibb/OpenProblemsInNLA/blob/main/matrix-functions-and-stability/MF-21/solution.md}{revised
proof} explicitly proves uniqueness of the continuous coefficient
functions on the logarithmic-squared bulk grid (Lemma 5, following
Proposition 4.2 of Barrera--Böttcher--Grudsky--Maximenko). This
identifies any hypothetical all-index expansion with the constructed
one. Equation (25) supplies only an upper bound; the obstruction uses
the separate inverse-trace contradiction in equations (26)--(33).
Corollary 6 now gives a matching lower bound of order \(h^{2m}\) for the
maximum error over a fixed finite set of low indices.

The revision also credits the simple-loop method explicitly and adds
Bogoya--Grudsky (2025). The
\href{https://github.com/ajt60gaibb/OpenProblemsInNLA/blob/main/reviews/2026-09-19-mf21/README.md}{dated
audit} records the mathematical checks, source comparison, and
coefficient-uniqueness repair. The
\href{https://github.com/ajt60gaibb/OpenProblemsInNLA/blob/main/reviews/2026-09-19-mf21/lean-audit.md}{initial
Lean audit} records the earlier partial development; it is superseded by
the completed verification above. The original problem, permanent ID,
grid, cutoff, and quantifiers below are unchanged.

\hypertarget{affirmative-resolution---12-september-2026-utc}{%
\subsection{Affirmative resolution - 12 September 2026
(UTC)}\label{affirmative-resolution---12-september-2026-utc}}

\textbf{Author:} George Stepaniants, Department of Computing and
Mathematical Sciences, California Institute of Technology, Pasadena,
California, USA.

\href{https://github.com/ajt60gaibb/OpenProblemsInNLA/blob/main/matrix-functions-and-stability/MF-21/solution.md}{Theorem
1 and Sections 2-5 of the complete proof} establish all three assertions
below for every integer \(m\ge3\), using the same smooth coefficient
functions with \(d_0=g_m\). The expansion is uniform over every
eigenvalue through order \(2m-1\), extends to order \(2m\) above the
stated logarithmic-squared index cutoff, and fails to extend uniformly
to all indices at that order. The matrix, grid normalization, parameter
quantifiers and cutoff are unchanged.

\href{https://github.com/ajt60gaibb/OpenProblemsInNLA/blob/main/matrix-functions-and-stability/MF-21/solution.pdf}{Proof
PDF} ·
\href{https://github.com/ajt60gaibb/OpenProblemsInNLA/blob/main/matrix-functions-and-stability/MF-21/solution.tex}{Standalone
proof TeX} ·
\href{https://github.com/ajt60gaibb/OpenProblemsInNLA/blob/main/references/stepaniants-mf21-2026-09-12/independent-review.md}{Independent
mathematical review} ·
\href{https://github.com/ajt60gaibb/OpenProblemsInNLA/blob/main/references/stepaniants-mf21-2026-09-12/README.md}{Submission
and public-source audit}.

The original proof passed a separate Codex-agent informal audit. It
derived the bulk expansion from an exact boundary determinant and the
final obstruction from a classical uniform inverse-kernel limit and a
trace identity; the revision above supplies a direct finite-matrix trace
calculation. Barrera, Böttcher, Grudsky, Maximenko and the cited later
authors retain credit for the conjecture and prior cases; Böttcher-Widom
and their cited predecessors retain credit for the inverse-kernel
theorem. Substantial AI assistance is disclosed. That initial audit was
informal; the completed Lean verification is documented above. External
human peer review is not claimed. The original statement, permanent ID,
references and dated history are retained below; the ratings are
historical.

\hypertarget{statement}{%
\subsection{Statement}\label{statement}}

For each integer \(m\ge3\), let

\[g_m(\theta)=\bigl(2\sin(\theta/2)\bigr)^{2m},\qquad
\widehat g_{m,k}=\frac1{2\pi}\int_{-\pi}^{\pi}g_m(\theta)e^{-ik\theta}\,d\theta.\]

Let \(T_n(g_m)=(\widehat g_{m,j-k})_{j,k=1}^n\), with eigenvalues
\(\lambda_{n,1}\le\cdots\le\lambda_{n,n}\) counted with multiplicity.
For real continuous functions \(d_0,\ldots,d_{2m}\) on \([0,\pi]\),
independent of \(n,j\), write

\[R_{p,n,j}=\lambda_{n,j}-\sum_{k=0}^p
\frac{d_k(j\pi/(n+2))}{(n+2)^k}.\]

\textbf{Conjecture 8.4 of Barrera--Böttcher--Grudsky--Maximenko.} Such
coefficient functions exist, with \(d_0=g_m\), for which all three
assertions hold:

\begin{enumerate}
\def\labelenumi{\arabic{enumi}.}
\item
  For every integer \(0\le p\le2m-1\), there are constants \(D_p>0\) and
  \(N_p\in\mathbb N\) such that

  \[|R_{p,n,j}|\le D_p(n+2)^{-p-1}\qquad(n\ge N_p,\ 1\le j\le n).\]
\item
  There are constants \(D_{2m}>0\) and \(N_{2m}\in\mathbb N\) such that

  \[|R_{2m,n,j}|\le D_{2m}(n+2)^{-2m-1}\]

  whenever \(n\ge N_{2m}\) and \(\lceil(\log(n+2))^2\rceil\le j\le n\).
\item
  No constants \(D>0,N\in\mathbb N\) make the bound
  \(|R_{2m,n,j}|\le D(n+2)^{-2m-1}\) valid for every \(n\ge N\) and
  every \(1\le j\le n\).
\end{enumerate}

All constants may depend on \(m\) and the expansion order; \(\log\)
denotes the natural logarithm. The same coefficient functions are used
throughout. The continuity and existence quantifiers make explicit the
regular-expansion convention of the source's Theorem 1.2 and §8; these
are not matrix-size-dependent fitted coefficients.

\hypertarget{numerical-significance}{%
\subsection{Numerical significance}\label{numerical-significance}}

Regular eigenvalue expansions support accurate computation without
forming large Toeplitz matrices. This conjecture identifies exactly
where extreme eigenvalues obstruct a uniform expansion for these
polynomial symbols. It retains all three parts as one problem.

\hypertarget{references-and-status-check}{%
\subsection{References and status
check}\label{references-and-status-check}}

\begin{itemize}
\tightlist
\item
  M. Barrera, A. Böttcher, S. M. Grudsky and E. A. Maximenko,
  \emph{Eigenvalues of even very nice Toeplitz matrices can be
  unexpectedly erratic}, Oper. Theory Adv. Appl. \textbf{268} (2018),
  51--77, \href{https://doi.org/10.1007/978-3-319-75996-8_2}{DOI};
  \href{https://arxiv.org/abs/1710.05243}{author preprint}, Conjecture
  8.4, p.~26, equation (8.4); Theorem 1.2 supplies the proved \(m=2\)
  analogue.
\item
  M. Barrera, S. Grudsky, V. Stukopin and I. Voronin, \emph{Asymptotics
  of the eigenvalues of seven-diagonal Toeplitz matrices of a special
  form}, Adv. Oper. Theory \textbf{9} (2024), 79,
  \href{https://doi.org/10.1007/s43036-024-00374-1}{DOI};
  \href{https://arxiv.org/abs/2111.07196}{preprint}, Theorems 2.3--2.6.
  This studies the sixth-order-zero case with more elaborate formulas.
\item
  M. Bogoya and S. Grudsky, \emph{Eigenvalues of non-Hermitian banded
  Toeplitz matrices approaching simple points of the limiting set},
  Comput. Math. Math. Phys. \textbf{65} (2025), 1453--1471,
  \href{https://doi.org/10.1134/S0965542525700745}{DOI};
  \href{https://www.math.cinvestav.mx/~grudsky/Papers/162.pdf}{author-hosted
  paper}, Theorems 2.1--2.2. This gives local phase equations and
  expansions near nondegenerate simple points; the revised proof records
  the methodological connection without a novelty claim.
\item
  A. Böttcher, \emph{Ten years with Sergei Grudsky in the eigenvalue
  bulk of Toeplitz matrices}, J. Math. Sci. \textbf{298} (2026),
  363--376, \href{https://doi.org/10.1007/s10958-025-07833-x}{DOI},
  section ``Beyond the simple-loop class,'' Theorems 2--5 and discussion
  of higher-order zeros.
\end{itemize}

On 2026-09-10, checked the original full preprint, the seven-diagonal
follow-up's full preprint and publication record, the 2026 survey, and
targeted title/conjecture/2025--2026 searches. No resolution of all
parts for every \(m\ge3\) was located. Later local second-order
expansions and results confined to \(m=3\) do not establish this full
statement. This bounded search is not a proof of openness.

\hypertarget{audit-2026-09-10}{%
\subsection{Audit --- 2026-09-10}\label{audit-2026-09-10}}

Independently rechecked
\href{https://arxiv.org/pdf/1710.05243}{Conjecture 8.4},
\href{https://arxiv.org/pdf/2111.07196}{the seven-diagonal paper's
Theorems 2.3--2.6}, and
\href{https://doi.org/10.1007/s10958-025-07833-x}{Böttcher's survey}.
Corrected the status to Open: the proved \(m=2\) analogue lies outside
the displayed range, and the \(m=3\) local and second-order formulas do
not establish the three-part threshold assertion. The
\href{https://doi.org/10.1007/s10958-026-08227-3}{March 2026
eigenvalue-superposition paper} assumes simple-loop symbols, excluding
these higher-order zeros. Title and conjecture searches found no
complete resolution. Challenging difficulty and specialist importance
are retained.

\hypertarget{resolution-audit---2026-09-12}{%
\subsection{Resolution audit -
2026-09-12}\label{resolution-audit---2026-09-12}}

The independent audit checked the complete all-\(m\) argument, including
the coalescing characteristic roots, normalized determinant remainder,
upper-endpoint cancellation, eigenvalue indexing, all three expansion
orders, and dominated trace limit. The primary-source and public
branch/fork/PR checks found no existing full resolution at the recorded
time; the linked submission record gives the precise scope and
limitations of that search.
\end{document}
